%! AP Statistics — Unit 9, Lesson 9.4 — Exit Ticket (Answer Key)
\documentclass[10pt]{article}
\usepackage{apstats-article}
\usepackage{apstats-key}

\begin{document}

\pageheader{Unit 9, Lesson 9.4}{Exit Ticket --- Answer Key}
\namedateperiod

\vspace{0.1in}

A school researcher randomly selects 35 students and records the average number of
hours per week they spend on extracurricular activities ($x$) and their GPA ($y$).
She wants to test whether extracurricular time is linearly related to GPA.

\begin{enumerate}

  \item Name the appropriate inference procedure and justify your choice.

        \ans{A $t$-test for slope, because we have one quantitative explanatory
        variable ($x$ = extracurricular hours) and one quantitative response variable
        ($y$ = GPA) and want to test whether the true slope of the population
        regression line differs from zero.}

  \item State H$_0$ and H$_a$ in terms of $\beta$. Define $\beta$ in context.

        H$_0$: \ans{$\beta = 0$; there is no linear relationship between extracurricular
        time and GPA in the population.}

        \vspace{0.05in}

        H$_a$: \ans{$\beta \ne 0$; there is a linear relationship between extracurricular
        time and GPA in the population.}

        \vspace{0.05in}

        $\beta = $ \ans{the true slope of the population regression line relating GPA
        to hours of extracurricular activities per week for students at this school}

  \item Which LINER conditions are met?

        \ans{L: Met --- residual plot shows random scatter (no curve).
        I: Met --- random sample; $n = 35 \le 10\% \times 350 = 35$, so 10\% condition
        is met (borderline).
        N: Met --- residuals approximately normal.
        E: Met --- residual plot shows roughly equal spread.
        R: Met --- random sample stated.}

\end{enumerate}

\begin{teachernote}
For condition I, $n = 35$ and $N > 350$ gives exactly $35 \le 0.10 \times 350 = 35$.
This is borderline. Accept student answers that acknowledge it is met (barely) or that
flag it as marginal. The key skill is checking the condition, not the exact conclusion.
\end{teachernote}

\end{document}
